\scanpage{34}
% archive foliation: 33
\uline{Exemples}

a) foncteurs $\underline{\operatorname{Hom}}$, $\underline{\otimes}$, et leurs dérivés. [2 arguments, $X=Y$] \unsure{on peut regarder un argument}.

$*$ \textcircled{d} Si $F$ et $G$ sont des faisceaux, on peut \unsure{former} l'ens \inserted{$N(F,G)$} des noyaux \inserted{et $P(F,G)$} des conoyaux des hom de $F$ en $G$ [\unsure{en ce qui} $\operatorname{Im}$ ou $\operatorname{CoIm}$, c'est kif kif, mais ce dernier est contenu ds un itéré des précédents]. (2 arg, $X=Y$)

\margin{\figure{\begin{tikzpicture}[x=.45cm,y=.45cm]
\draw[->] (1,5.4) .. controls (-1,1.8) and (-.8,.9) .. (1,0);
\draw[->] (.6,-4.1) .. controls (-.5,-3.0) and (.0,-2.0) .. (1.7,-1.0);
\draw[->] (-.4,.3) .. controls (-.5,1.6) and (-.3,1.9) .. (1.4,1.9);
\end{tikzpicture}}}

b) Foncteur image inverse (1 arg, $f:X\to Y$ donné), et foncteurs dérivés

c) Foncteur image directe (1 arg. $f:X\to Y$ \unsure{propre} donné), et foncteurs dérivés

$*$ \textcircled{e} Ens. des extensions de $F$ par $G$ (2 arg, $X=Y$)

\textcircled{g} Ens des quotients associés à la filtration primaire croissante, resp. décroissante, de $F$ (1 arg, $X=Y$)

$*$ \textcircled{f} \deleted{Ens \illegible{2}} \inserted{Foncteur} $\underline O_{(\operatorname{Supp}F)_{\text{réd}}}$ ($K$ alg. clos, 1 arg, $X=Y$)

\deleted{Ens des composantes}

\textcircled{h} Ens des \deleted{\unsure{quotients}} \deleted{faisceaux} composantes primaires de $F$ associés aux composantes de son support ($K$ alg. clos, $\unsure{\chi}=0$, $X=Y$)

\deleted{\textcircled{i} \illegible{8} $F$.}
